% Full title: Chinese Translation of Essential Spectrum for Two Different Periodic Ends
% Purpose:
%   Provide the Chinese counterpart of the complete essential-spectrum proof
%   for a scalar strip with two different periodic ends.
% Main workflow:
%   Establish local compactness and cutoff commutator bounds, derive a Zhislin
%   criterion, construct one-sided periodic Weyl sequences, and prove both
%   inclusions in the essential-spectrum identity.
% Based on:
%   The Zhislin-sequence mechanism of De Nittis et al. (2021) and the periodic
%   Floquet--Bloch characterization of Fliss (2013).
% Main changes:
%   Adapts the one-dimensional first-order argument to the present
%   two-dimensional second-order weighted strip operator with exact periodic
%   tails and possibly different transverse periods.
% Mathematical goal:
%   Prove Theorem~\ref{thm:two-ended-essential-spectrum} without assuming the
%   essential-spectrum union formula.

\section{两个不同周期端对应的本质谱}
\label{app:essential-spectrum}

在本附录中，固定
$\beta\in[-\pi/d,\pi/d)$，并简记
\[
  A:=A(\beta),\qquad A^\star:=A^\star(\beta),
  \qquad \star\in\{-,+\}.
\]
这些算子分别作用于
$H=L^2(B,\rho\,\dd x\dd y)$ 和
$H^\star=L^2(B,\rho^\star\,\dd x\dd y)$。定义
\begin{equation}
  H_0:=H^-\oplus H^+,
  \qquad A_0:=A^-\oplus A^+.
  \label{eq:appendix-free-operator}
\end{equation}
以 $Q$ 表示条带空间上的乘法算子 $u(x,y)\mapsto xu(x,y)$，并令
$Q_0:=Q\oplus Q$ 作用于 $H_0$。取充分大的 $R_0>0$，使得
\[
  \overline{\mathcal C^0}\subset(-R_0,R_0)\times[0,d],
  \qquad
  \rho=\rho^-\quad\text{在 }x<-R_0\text{ 上},
  \qquad
  \rho=\rho^+\quad\text{在 }x>R_0\text{ 上}.
\]
后两个等式由两个周期半波导的构造保证。由于各系数具有公共的正下界和有限
上界，下文所有加权范数均与相应的非加权范数等价。对于标量算子 $T$，当
$T=A$ 时记 $\rho_T=\rho$，当 $T=A^\star$ 时记
$\rho_T=\rho^\star$；涉及 $A_0$ 的公式均逐分量理解。

\begin{lemma}[局部紧性与截断交换子]
\label{lem:appendix-local-compactness}
设 $T$ 为 $A,A^-,A^+$ 或 $A_0$ 中任一算子，作用于相应的 Hilbert 空间
$H_T$；以 $Q_T$ 表示相应的位置算子。则：
\begin{enumerate}[label=\textup{(\roman*)},leftmargin=*]
  \item 对每个有界区间 $K\subset\R$，
  \[
    \mathbf 1_K(Q_T)(T+\ii)^{-1}:H_T\longrightarrow H_T
    \quad\text{是紧算子};
  \]
  \item 若 $\zeta\in C_c^\infty(\R;[0,1])$ 在 $[-1,1]$ 上等于 $1$，
  在 $[-2,2]$ 外等于 $0$，并令
  $\zeta_R(x):=\zeta(x/R)$，则
  \begin{equation}
    \bigl\|[T,\zeta_R(Q_T)](T+\ii)^{-1}\bigr\|_{\mathcal B(H_T)}
    \longrightarrow0
    \qquad(R\to\infty).
    \label{eq:appendix-commutator-limit}
  \end{equation}
\end{enumerate}
\end{lemma}
\begin{proof}
先考虑一个标量分量。其闭型为
\[
  a_T(u,v)=\int_B\nabla u\cdot\nabla\overline v\,\dd x\dd y,
  \qquad D(a_T)=H^1_\beta(B).
\]
若 $u=(T+\ii)^{-1}f$，则由非负性和变分恒等式得到
\[
  \|u\|_{H_T}\leq\|f\|_{H_T},
  \qquad
  \|\nabla u\|_{L^2(B)}^2
  =\operatorname{Re}\langle f,u\rangle_{H_T}
  \leq\|f\|_{H_T}^2.
\]
所以预解算子有界地映入 $H^1_\beta(B)$。先限制到
$K\times(0,d)$，再使用 Rellich 紧嵌入
$H^1(K\times(0,d))\hookrightarrow L^2(K\times(0,d))$，即可得到紧性。
系数上下界保证局部加权与非加权 $L^2$ 范数等价，故 (i) 成立。直和情形逐
分量即得。

只依赖于 $x$ 的光滑函数乘法保持算子定义域不变：若
$u\in H^1(\Delta,B)$，则
\[
  \Delta(\eta u)=\eta\Delta u+2\eta'\partial_xu+\eta''u\in L^2(B),
\]
并且水平方向的 $\beta$-准周期迹条件不变。因此，在 $D(T)$ 上，
\begin{equation}
  [T,\zeta_R]u
  =-\frac1{\rho_T}
    \left(2\zeta_R'\partial_xu+\zeta_R''u\right),
  \label{eq:appendix-commutator-formula}
\end{equation}
其中对 $A_0$ 逐分量理解。由于
$\|\zeta_R'\|_\infty=O(R^{-1})$ 且
$\|\zeta_R''\|_\infty=O(R^{-2})$，前述预解式估计给出
\[
  \bigl\|[T,\zeta_R](T+\ii)^{-1}f\bigr\|_{H_T}
  \leq C\left(R^{-1}+R^{-2}\right)\|f\|_{H_T}.
\]
这就证明了 (ii)。
\end{proof}

\begin{lemma}[条带上的 Zhislin 判据]
\label{lem:appendix-Zhislin}
对于 \Cref{lem:appendix-local-compactness} 中任一算子 $T$ 和任意
$\lambda\in\R$，下列两条等价：
\begin{enumerate}[label=\textup{(\roman*)},leftmargin=*]
  \item $\lambda\in\sigma_{\mathrm{ess}}(T)$；
  \item 存在 $u_m\in D(T)$，使得
  \begin{equation}
    \|u_m\|_{H_T}=1,\qquad
    \mathbf 1_{[-m,m]}(Q_T)u_m=0,\qquad
    \|(T-\lambda)u_m\|_{H_T}\longrightarrow0.
    \label{eq:appendix-Zhislin-sequence}
  \end{equation}
\end{enumerate}
当 $T=A_0$ 时，支撑条件同时施加于两个分量。
\end{lemma}
\begin{proof}
满足 (ii) 的序列弱收敛到零：先对紧支撑的 $L^2$ 函数验证这一点，再由稠密性
推广到任意函数即可。因此它是一个奇异 Weyl 序列，从而推出 (i)。

反过来，设 $v_j$ 是 $T$ 在 $\lambda$ 处的奇异 Weyl 序列。对每个固定的
$R$，局部紧性给出
\[
  \zeta_R(Q_T)v_j
  =\zeta_R(Q_T)(T+\ii)^{-1}(T+\ii)v_j
  \longrightarrow0
  \quad(j\to\infty)
\]
的强收敛。事实上，
$(T+\ii)v_j=(T-\lambda)v_j+(\lambda+\ii)v_j$ 弱收敛到零且有界。
选取一个子序列，仍记为 $v_m$，使得
\[
  \|\zeta_m(Q_T)v_m\|_{H_T}\leq m^{-1},\qquad
  \|(T-\lambda)v_m\|_{H_T}\leq m^{-1}.
\]
令 $w_m=(1-\zeta_m(Q_T))v_m$。于是 $w_m$ 在 $|x|\leq m$ 上为零，并且
$\|w_m\|_{H_T}\to1$。此外，
\[
  (T-\lambda)w_m
  =(1-\zeta_m)(T-\lambda)v_m-[T,\zeta_m]v_m.
\]
第一项趋于零。对于第二项，写成
\[
  [T,\zeta_m]v_m
  =[T,\zeta_m](T+\ii)^{-1}(T+\ii)v_m.
\]
第一个因子由 \eqref{eq:appendix-commutator-limit} 趋于零，第二个因子保持
有界。将 $w_m$ 归一化即得 \eqref{eq:appendix-Zhislin-sequence}。这正是
\cite[Lemmas~3.11--3.12]{DeNittisEtAl2021} 中 Zhislin 构造在当前条带上的
对应形式。
\end{proof}

\begin{lemma}[置于指定端部的周期 Weyl 序列]
\label{lem:appendix-one-sided-Weyl}
设 $\star\in\{-,+\}$ 且
$\lambda\in\sigma(A^\star)$。则存在归一化的
$v_m^\star\in D(A^\star)$，使得
\[
  \|(A^\star-\lambda)v_m^\star\|_{H^\star}\longrightarrow0
\]
以及
\[
  \operatorname{supp}v_m^+\subset[m,\infty)\times[0,d],
  \qquad
  \operatorname{supp}v_m^-\subset(-\infty,-m]\times[0,d].
\]
\end{lemma}
\begin{proof}
只写出右端构造。由谱的 Weyl 判据，可取 $f_m\in D(A^+)$，满足
\[
  \|f_m\|_{H^+}=1,\qquad
  \|(A^+-\lambda)f_m\|_{H^+}\leq m^{-1}.
\]
因为 $A^+$ 非负且由梯度型生成，
\[
  \|\nabla f_m\|_{L^2(B)}^2
  =\langle A^+f_m,f_m\rangle_{H^+}
  \leq\lambda+m^{-1}
\]
对所有充分大的 $m$ 成立。周期平移
\[
  (\tau_q^+f)(x,y):=f(x-qL^+,y),\qquad q\in\Z,
\]
在 $H^+$ 中是酉算子，并与 $A^+$ 交换。取充分大的 $q_m$，使得
\[
  \bigl\|\mathbf 1_{(-\infty,2m]}(Q)\tau_{q_m}^+f_m
  \bigr\|_{H^+}\leq m^{-1}.
\]
取 $\eta\in C^\infty(\R;[0,1])$，满足
$\eta(t)=0$（$t\leq0$）以及 $\eta(t)=1$（$t\geq1$），并令
$\eta_m^+(x):=\eta((x-m)/m)$。则
\[
  g_m:=\eta_m^+\tau_{q_m}^+f_m
\]
支撑于 $[m,\infty)\times[0,d]$，且 $\|g_m\|_{H^+}\to1$。
\Cref{lem:appendix-local-compactness} 中保持定义域的计算说明
$g_m\in D(A^+)$，并且
\[
  (A^+-\lambda)g_m
  =\eta_m^+\tau_{q_m}^+(A^+-\lambda)f_m
   +[A^+,\eta_m^+]\tau_{q_m}^+f_m.
\]
第一项趋于零。由 \eqref{eq:appendix-commutator-formula}、一致的
$H^1$ 界以及
$\|\partial_x^r\eta_m^+\|_\infty=O(m^{-r})$（$r=1,2$），第二项也趋于零。
归一化 $g_m$ 即得到 $v_m^+$。沿负方向平移 $L^-$ 的整数倍并使用
$\eta_m^-(x):=\eta((-x-m)/m)$，即可得到左端序列。这是
\cite[Proposition~3.10]{DeNittisEtAl2021} 证明中的平移--截断构造在二阶
条带算子上的对应形式。
\end{proof}

\begin{proof}[\Cref{thm:two-ended-essential-spectrum} 的证明]
分别对两个周期延拓使用 Floquet--Bloch 分解，并利用
\eqref{eq:background-bands}，可得
\begin{equation}
  \sigma_{\mathrm{ess}}(A_0)
  =\sigma_{\mathrm{ess}}(A^-)
   \cup\sigma_{\mathrm{ess}}(A^+)
  =\sigma(A^-)\cup\sigma(A^+)
  =\Sigma_-(\beta)\cup\Sigma_+(\beta).
  \label{eq:appendix-free-essential-spectrum}
\end{equation}
这里第一个等式是有限正交直和的本质谱公式；第二个等式来自
\cite[Proposition~3.1, especially equations~(3.3)--(3.4)]{Fliss2013}
的周期谱结论。

先证
$\sigma_{\mathrm{ess}}(A)\subset\sigma_{\mathrm{ess}}(A_0)$。
取 $j_-,j_+\in C^\infty(\R;[0,1])$，满足
\[
\begin{array}{lll}
  j_-(x)=1&\text{当 }x\leq-2R_0,&
  j_-(x)=0\text{ 当 }x\geq-R_0,\\
  j_+(x)=0&\text{当 }x\leq R_0,&
  j_+(x)=1\text{ 当 }x\geq2R_0.
\end{array}
\]
设 $\lambda\in\sigma_{\mathrm{ess}}(A)$，并取
\Cref{lem:appendix-Zhislin} 给出的 Zhislin 序列 $u_m$。当
$m>2R_0$ 时，
\[
  u_m=j_-u_m+j_+u_m,
\]
且右端两项支撑不交。定义
\[
  u_m^0:=c_m(j_-u_m,j_+u_m)\in D(A_0),
  \qquad
  c_m:=\|(j_-u_m,j_+u_m)\|_{H_0}^{-1}.
\]
在 $j_-u_m$ 的支撑上，权重 $\rho$ 与 $\rho^-$ 相等；在
$j_+u_m$ 的支撑上，权重 $\rho$ 与 $\rho^+$ 相等。因此，对所有充分大的
$m$，均有 $c_m=1$。函数 $j_\pm$ 的导数支撑于固定集合
$[-2R_0,2R_0]$，而 $u_m$ 在该集合上为零。所以交换子项为零，并且
\[
\begin{aligned}
  \|(A_0-\lambda)u_m^0\|_{H_0}^2
  &=\|(A^--\lambda)j_-u_m\|_{H^-}^2
    +\|(A^+-\lambda)j_+u_m\|_{H^+}^2\\
  &\leq C\|(A-\lambda)u_m\|_H^2\longrightarrow0.
\end{aligned}
\]
序列 $u_m^0$ 在 $|x|\leq m$ 上仍为零。由
\Cref{lem:appendix-Zhislin}，
$\lambda\in\sigma_{\mathrm{ess}}(A_0)$。

对于反向包含，设 $\lambda\in\sigma(A^+)$。由
\Cref{lem:appendix-one-sided-Weyl}，存在归一化近似特征向量序列
$v_m^+$，其支撑包含于 $[m,\infty)\times[0,d]$。一旦 $m>R_0$，两个
权重以及两个微分实现就在该支撑上相同。因此
\[
  \|v_m^+\|_H=1,\qquad
  \|(A-\lambda)v_m^+\|_H
  =\|(A^+-\lambda)v_m^+\|_{H^+}\longrightarrow0.
\]
因为这些支撑向右端逃逸，所以 $v_m^+\rightharpoonup0$ 于 $H$。它是
$A$ 的一个奇异 Weyl 序列，故
$\lambda\in\sigma_{\mathrm{ess}}(A)$。同理，左端构造给出
$\sigma(A^-)\subset\sigma_{\mathrm{ess}}(A)$。所以
\[
  \sigma(A^-)\cup\sigma(A^+)
  \subset\sigma_{\mathrm{ess}}(A)
  \subset\sigma_{\mathrm{ess}}(A_0).
\]
结合 \eqref{eq:appendix-free-essential-spectrum} 即证
\eqref{eq:projected-spectrum}。
\end{proof}
